\section{AHSME 1993}
        \begin{enumerate}
        	\item (\textit{AHSME 1993}) For integers $a, b$ and $c$, define $\boxed{a,b,c}$ to mean $a^b-b^c+c^a$. Then $\boxed{1,-1,2}$ equals


 $\text{(A)} \ -4 \qquad \text{(B)} \ -2 \qquad \text{(C)} \ 0 \qquad \text{(D)} \ 2 \qquad \text{(E)} \ 4$


 

	\item (\textit{AHSME 1993}) In $\triangle ABC$, $\angle A=55^\circ$, $\angle C=75^\circ$, $D$ is on side $\overline{AB}$ and $E$ is on side $\overline{BC}$ If $DB=BE$, then $\angle BED=$


 $\text{(A)}\ 50^\circ \qquad \text{(B)}\ 55^\circ \qquad \text{(C)}\ 60^\circ \qquad \text{(D)}\ 65^\circ \qquad \text{(E)}\ 70^\circ$


 

	\item (\textit{AHSME 1993}) $\frac{15^{30}}{45^{15}} =$


 $\text{(A) } \left(\frac{1}{3}\right)^{15}\quad \text{(B) } \left(\frac{1}{3}\right)^{2}\quad \text{(C) } 1\quad \text{(D) } 3^{15}\quad \text{(E) } 5^{15}$


 

	\item (\textit{AHSME 1993}) Define the operation "$\circ$" by $x\circ y=4x-3y+xy$, for all real numbers $x$ and $y$. For how many real numbers $y$ does $3\circ y=12$?


 $\text{(A) } 0\quad \text{(B) } 1\quad \text{(C) } 3\quad \text{(D) } 4\quad \text{(E) more than 4}$


 

	\item (\textit{AHSME 1993}) Last year a bicycle cost $160 and a cycling helmet $40. This year the cost of the bicycle increased by $5\%$, and the cost of the helmet increased by $10\%$. The percent increase in the combined cost of the bicycle and the helmet is:


 $\text{(A) } 6\%\quad \text{(B) } 7\%\quad \text{(C) } 7.5\%\quad \text{(D) } 8\%\quad \text{(E) } 15\%$


 

	\item (\textit{AHSME 1993}) $\sqrt{\frac{8^{10}+4^{10}}{8^4+4^{11}}}=$


 $\text{(A) } \sqrt{2}\quad \text{(B) } 16\quad \text{(C) } 32\quad \text{(D) } (12)^{\frac{2}{3}}\quad \text{(E) } 512.5$


 

	\item (\textit{AHSME 1993}) The symbol $R_k$ stands for an integer whose base-ten representation is a sequence of $k$ ones. For example, $R_3=111,R_5=11111$, etc. When $R_{24}$ is divided by $R_4$, the quotient $Q=R_{24}/R_4$ is an integer whose base-ten representation is a sequence containing only ones and zeroes. The number of zeros in $Q$ is:


 $\text{(A) } 10\quad \text{(B) } 11\quad \text{(C) } 12\quad \text{(D) } 13\quad \text{(E) } 15$


 

	\item (\textit{AHSME 1993}) Let $C_1$ and $C_2$ be circles of radius 1 that are in the same plane and tangent to each other. How many circles of radius 3 are in this plane and tangent to both $C_1$ and $C_2$?


 $\text{(A) } 2\quad \text{(B) } 4\quad \text{(C) } 5\quad \text{(D) } 6\quad \text{(E) } 8$


 

	\item (\textit{AHSME 1993}) Country $A$ has $c\%$ of the world's population and $d\%$ of the worlds wealth. Country $B$ has $e\%$ of the world's population and $f\%$ of its wealth. Assume that the citizens of $A$ share the wealth of $A$ equally,and assume that those of $B$ share the wealth of $B$ equally. Find the ratio of the wealth of a citizen of $A$ to the wealth of a citizen of $B$.


 $\text{(A) } \frac{cd}{ef}\quad \text{(B) } \frac{ce}{ef}\quad \text{(C) } \frac{cf}{de}\quad \text{(D) } \frac{de}{cf}\quad \text{(E) } \frac{df}{ce}$


 

	\item (\textit{AHSME 1993}) Let $r$ be the number that results when both the base and the exponent of $a^b$ are tripled, where $a,b>0$. If $r$ equals the product of $a^b$ and $x^b$ where $x>0$, then $x=$


 $\text{(A) } 3\quad \text{(B) } 3a^2\quad \text{(C) } 27a^2\quad \text{(D) } 2a^{3b}\quad \text{(E) } 3a^{2b}$


 

	\item (\textit{AHSME 1993}) If $\log_2(\log_2(\log_2(x)))=2$, then how many digits are in the base-ten representation for $x$?


 $\text{(A) } 5\quad \text{(B) } 7\quad \text{(C) } 9\quad \text{(D) } 11\quad \text{(E) } 13$


 

	\item (\textit{AHSME 1993}) If $f(2x)=\frac{2}{2+x}$ for all $x>0$, then $2f(x)=$


 $\text{(A) } \frac{2}{1+x}\quad \text{(B) } \frac{2}{2+x}\quad \text{(C) } \frac{4}{1+x}\quad \text{(D) } \frac{4}{2+x}\quad \text{(E) } \frac{8}{4+x}$


 

	\item (\textit{AHSME 1993}) A square of perimeter 20 is inscribed in a square of perimeter 28. What is the greatest distance between a vertex of the inner square and a vertex of the outer square?


 $\text{(A) } \sqrt{58}\quad \text{(B) } \frac{7\sqrt{5}}{2}\quad \text{(C) } 8\quad \text{(D) } \sqrt{65}\quad \text{(E) } 5\sqrt{3}$


 

	\item (\textit{AHSME 1993}) 
                    \begin{figure}[H]
                        \centering
                        \includegraphics[width=0.5\textwidth]{Problem/AHSME_I_1993_14_Problem_Figure_1.png}
                    \end{figure}
                    


 The convex pentagon $ABCDE$ has $\angle{A}=\angle{B}=120^\circ,EA=AB=BC=2$ and $CD=DE=4$. What is the area of ABCDE?


 $\text{(A) } 10\quad \text{(B) } 7\sqrt{3}\quad \text{(C) } 15\quad \text{(D) } 9\sqrt{3}\quad \text{(E) } 12\sqrt{5}$


 

	\item (\textit{AHSME 1993}) For how many values of $n$ will an $n$-sided regular polygon have interior angles with integral measures?


 $\text{(A) } 16\quad \text{(B) } 18\quad \text{(C) } 20\quad \text{(D) } 22\quad \text{(E) } 24$


 

	\item (\textit{AHSME 1993}) Consider the non-decreasing sequence of positive integers

 \[1,2,2,3,3,3,4,4,4,4,5,5,5,5,5,\cdots\]
in which the $n^{th}$ positive integer appears $n$ times. The remainder when the $1993^{rd}$ term is divided by $5$ is


 $\text{(A) } 0\quad \text{(B) } 1\quad \text{(C) } 2\quad \text{(D) } 3\quad \text{(E) } 4$


 

	\item (\textit{AHSME 1993}) 
                    \begin{figure}[H]
                        \centering
                        \includegraphics[width=0.5\textwidth]{Problem/AHSME_I_1993_17_Problem_Figure_1.png}
                    \end{figure}
                    


 Amy painted a dartboard over a square clock face using the "hour positions" as boundaries.[See figure.] If $t$ is the area of one of the eight triangular regions such as that between 12 o'clock and 1 o'clock, and $q$ is the area of one of the four corner quadrilaterals such as that between 1 o'clock and 2 o'clock, then $\frac{q}{t}=$


 $\text{(A) } 2\sqrt{3}-2\quad \text{(B) } \frac{3}{2}\quad \text{(C) } \frac{\sqrt{5}+1}{2}\quad \text{(D) } \sqrt{3}\quad \text{(E) } 2$


 

	\item (\textit{AHSME 1993}) Al and Barb start their new jobs on the same day. Al's schedule is 3 work-days followed by 1 rest-day. Barb's schedule is 7 work-days followed by 3 rest-days. On how many of their first 1000 days do both have rest-days on the same day?


 $\text{(A) } 48\quad \text{(B) } 50\quad \text{(C) } 72\quad \text{(D) } 75\quad \text{(E) } 100$


 

	\item (\textit{AHSME 1993}) How many ordered pairs $(m,n)$ of positive integers are solutions to

 \[\frac{4}{m}+\frac{2}{n}=1?\]


 $\text{(A) } 1\quad \text{(B) } 2\quad \text{(C) } 3\quad \text{(D) } 4\quad \text{(E) } \text{more than }6$


 

	\item (\textit{AHSME 1993}) Consider the equation $10z^2-3iz-k=0$, where $z$ is a complex variable and $i^2=-1$. Which of the following statements is true?


 $\text{(A) For all positive real numbers k, both roots are pure imaginary} \quad\\ \text{(B) For all negative real numbers k, both roots are pure imaginary} \quad\\ \text{(C) For all pure imaginary numbers k, both roots are real and rational} \quad\\ \text{(D) For all pure imaginary numbers k, both roots are real and irrational} \quad\\ \text{(E) For all complex numbers k, neither root is real}$


 

	\item (\textit{AHSME 1993}) Let $a_1,a_2,\cdots,a_k$ be a finite arithmetic sequence with $a_4 +a_7+a_{10} = 17$ and $a_4+a_5+\cdots+a_{13} +a_{14} = 77$. 


 If $a_k = 13$, then $k =$


 $\text{(A) } 16\quad \text{(B) } 18\quad \text{(C) } 20\quad \text{(D) } 22\quad \text{(E) } 24$


 

	\item (\textit{AHSME 1993}) 
                    \begin{figure}[H]
                        \centering
                        \includegraphics[width=0.5\textwidth]{Problem/AHSME_I_1993_22_Problem_Figure_1.png}
                    \end{figure}
                    


 Twenty cubical blocks are arranged as shown. First, 10 are arranged in a triangular pattern; then a layer of 6, arranged in a triangular pattern, is centered on the 10; then a layer of 3, arranged in a triangular pattern, is centered on the 6; and finally one block is centered on top of the third layer. The blocks in the bottom layer are numbered 1 through 10 in some order. Each block in layers 2,3 and 4 is assigned the number which is the sum of numbers assigned to the three blocks on which it rests. Find the smallest possible number which could be assigned to the top block.


 $\text{(A) } 55\quad \text{(B) } 83\quad \text{(C) } 114\quad \text{(D) } 137\quad \text{(E) } 144$


 

	\item (\textit{AHSME 1993}) 
                    \begin{figure}[H]
                        \centering
                        \includegraphics[width=0.5\textwidth]{Problem/AHSME_I_1993_23_Problem_Figure_1.png}
                    \end{figure}
                    


 Points $A,B,C$ and $D$ are on a circle of diameter $1$, and $X$ is on diameter $\overline{AD}.$


 If $BX=CX$ and $3\angle{BAC}=\angle{BXC}=36^\circ$, then $AX=$


 

$\text{(A) } \cos(6^\circ)\cos(12^\circ)\sec(18^\circ)\quad\\ \text{(B) } \cos(6^\circ)\sin(12^\circ)\csc(18^\circ)\quad\\ \text{(C) } \cos(6^\circ)\sin(12^\circ)\sec(18^\circ)\quad\\ \text{(D) } \sin(6^\circ)\sin(12^\circ)\csc(18^\circ)\quad\\ \text{(E) } \sin(6^\circ)\sin(12^\circ)\sec(18^\circ)$


 

	\item (\textit{AHSME 1993}) A box contains $3$ shiny pennies and $4$ dull pennies. One by one, pennies are drawn at random from the box and not replaced. If the probability is $a/b$ that it will take more than four draws until the third shiny penny appears and $a/b$ is in lowest terms, then $a+b=$


 $\text{(A) } 11\quad \text{(B) } 20\quad \text{(C) } 35\quad \text{(D) } 58\quad \text{(E) } 66$


 

	\item (\textit{AHSME 1993}) 
                    \begin{figure}[H]
                        \centering
                        \includegraphics[width=0.5\textwidth]{Problem/AHSME_I_1993_25_Problem_Figure_1.png}
                    \end{figure}
                    


 Let $S$ be the set of points on the rays forming the sides of a $120^{\circ}$ angle, and let $P$ be a fixed point inside the angle 
on the angle bisector. Consider all distinct equilateral triangles $PQR$ with $Q$ and $R$ in $S$. 
(Points $Q$ and $R$ may be on the same ray, and switching the names of $Q$ and $R$ does not create a distinct triangle.) 
There are


 $\text{(A) exactly 2 such triangles} \quad\\ \text{(B) exactly 3 such triangles} \quad\\ \text{(C) exactly 7 such triangles} \quad\\ \text{(D) exactly 15 such triangles} \quad\\ \text{(E) more than 15 such triangles}$


 

	\item (\textit{AHSME 1993}) Find the largest positive value attained by the function

 \[f(x)=\sqrt{8x-x^2}-\sqrt{14x-x^2-48} ,\quad x  \text{  a real number}\]


 $\text{(A) } \sqrt{7}-1\quad \text{(B) } 3\quad \text{(C) } 2\sqrt{3}\quad \text{(D) } 4\quad \text{(E) } \sqrt{55}-\sqrt{5}$


 

	\item (\textit{AHSME 1993}) 
                    \begin{figure}[H]
                        \centering
                        \includegraphics[width=0.5\textwidth]{Problem/AHSME_I_1993_27_Problem_Figure_1.png}
                    \end{figure}
                    
The sides of $\triangle ABC$ have lengths $6,8,$ and $10$. A circle with center $P$ and radius $1$ rolls around the inside of $\triangle ABC$, always remaining tangent to at least one side of the triangle. When $P$ first returns to its original position, through what distance has $P$ traveled?


 $\text{(A) } 10\quad \text{(B) } 12\quad \text{(C) } 14\quad \text{(D) } 15\quad \text{(E) } 17$


 

	\item (\textit{AHSME 1993}) How many triangles with positive area are there whose vertices are points in the $xy$-plane whose coordinates are integers $(x,y)$ satisfying $1\le x\le 4$ and $1\le y\le 4$?


 $\text{(A) } 496\quad \text{(B) } 500\quad \text{(C) } 512\quad \text{(D) } 516\quad \text{(E) } 560$


 

	\item (\textit{AHSME 1993}) Which of the following could NOT be the lengths of the external diagonals of a right regular prism [a "box"]? (An $\textit{external diagonal}$ is a diagonal of one of the rectangular faces of the box.)


 $\text{(A) }\{4,5,6\} \quad \text{(B) } \{4,5,7\} \quad \text{(C) } \{4,6,7\} \quad \text{(D) } \{5,6,7\} \quad \text{(E) } \{5,7,8\}$


 

	\item (\textit{AHSME 1993}) Given $0\le x_0<1$, let 

 \[x_n=\left\{ \begin{array}{ll} 2x_{n-1} &\text{ if }2x_{n-1}<1 \\ 2x_{n-1}-1 &\text{ if }2x_{n-1}\ge 1 \end{array}\right.\]
for all integers $n>0$. For how many $x_0$ is it true that $x_0=x_5$?


 $\text{(A) 0} \quad \text{(B) 1} \quad \text{(C) 5} \quad \text{(D) 31} \quad \text{(E) }\infty$


 

\end{enumerate}